
\chapter{further prescriptive analytics to control the \ac{ESP}}
\subsection{Introduction}
The operation of the ESP consists of adjusting the pump rotational speed and the opening of the production valve (choke) at the wellhead that feeds the manifold. Moreover, the safe and stable operation of ESP-lifted oil wells is carried out by the so-called ESP operating envelope-like set(Takacs, 2009), which comprises time-variant constraints (upthrust and downthrust) that rely on the phase envelope-like set(Takacs, 2009), which comprises time-variant constraints (upthrust and downthrust) that rely on the p An automatic control of Electric Submersible Pumps (ESPs), employed for artificial lift oil production, is desired in order to ensure safe operation while satisfying pump and well constraints. 

Currently pumps are employed with variable speed drive that change frequency. The objective  is  to control the inlet pressure, at a desired  reference value    and    to    minimize    the    pump    power    consumption, while    satisfying    the    pump    and    well    constraints. The pump frequency and valve opening were considered  as  manipulated  variables  in  order  to  achieve  this  objective. 
1- Models describe the whole system must be identified (Hybrid Data Driven and Mechanistic models)
2- Controller must be used [e.g. MPC]
Or may be an application of Reinforcement learning 
Study aims at maximizing oil production and minimize power consumption.

\subsection{system description}

The considered system consists of one oil well with an Electric Submersible Pump (ESP) and a production choke valve, as seen in Fig. 1. The ESP is used for artificial lift, with the purpose of reducing the hydro-static pressure from the fluid in the well and the production line, thus reducing the bottom hole pressure, increasing the inflow from the reservoir
and the production from the well.

A. Control targets
There are a number of variables that affect the life-time of ESP installations, such as power consumption, flow rate, pressure, temperature, thrust forces and vibration. Operation
outside of certain limits on these variables may lead to failure or reduced life-time of the ESP, which has a huge economic impact, both due to the costs of replacing the pump, and the
loss of production. Thus, the main control priority in this

system is to maintain good operating conditions for the ESP. This study focuses on keeping a constraint on the electric current, which is proportional to the power consumption
of the pump, and keeping constraints on the upthrust and downthrust forces acting on the ESP shaft.

The next priority is to ensure acceptable operational conditions for the well. This includes keeping the pump inlet pressure within certain constraints, and to avoid back-flow into the well by keeping the wellhead pressure higher than
the topside pressure.

Less prioritized, but still important control targets, are to
maximize the pump efficiency, and optimize certain production condition parameters. This includes keeping the inlet pressure at a certain set-point, and minimize the electric
current flow in the ESP motor. Another control target is to avoid high wellhead pressure relative to the topside pressure; a high pressure difference means there is a large power loss in the topside choke, which is a waste of power provided by the ESP. Choking the well may be necessary to keep the ESP within its operational limits, but is otherwise not desired.

B. Third order model A simple, third order model of an ESP lifted well is
developed by Statoil in [12], [13]. A similar model is



\section{Implementation of reinforcement learning}
\begin{preview}
    \begin{algorithm}[H]
        \begin{algorithmic}
        \Require
        \Statex States $\mathcal{X} = \{1, \dots, n_x\}$
        \Statex Actions $\mathcal{A} = \{1, \dots, n_a\},\qquad A: \mathcal{X} \Rightarrow \mathcal{A}$
        \Statex Reward function $R: \mathcal{X} \times \mathcal{A} \rightarrow \mathbb{R}$
        \Statex Black-box (probabilistic) transition function $T: \mathcal{X} \times \mathcal{A} \rightarrow \mathcal{X}$
        \Statex Learning rate $\alpha \in [0, 1]$, typically $\alpha = 0.1$
        \Statex Discounting factor $\gamma \in [0, 1]$
        \Procedure{QLearning}{$\mathcal{X}$, $A$, $R$, $T$, $\alpha$, $\gamma$}
            \State Initialize $Q: \mathcal{X} \times \mathcal{A} \rightarrow \mathbb{R}$ arbitrarily
            \While{$Q$ is not converged}
                \State Start in state $s \in \mathcal{X}$
                \While{$s$ is not terminal}
                    \State Calculate $\pi$ according to Q and exploration strategy (e.g. $\pi(x) \gets \argmax_{a} Q(x, a)$)
                    \State $a \gets \pi(s)$
                    \State $r \gets R(s, a)$ \Comment{Receive the reward}
                    \State $s' \gets T(s, a)$ \Comment{Receive the new state}
                    \State $Q(s', a) \gets (1 - \alpha) \cdot Q(s, a) + \alpha \cdot (r + \gamma \cdot \max_{a'} Q(s', a'))$
                    \State $s \gets s'$
                \EndWhile
            \EndWhile
            \Return $Q$
        \EndProcedure
        \end{algorithmic}
    \caption{$Q$-learning: Learn function $Q: \mathcal{X} \times \mathcal{A} \rightarrow \mathbb{R}$}
    \label{alg:q-learning}
    \end{algorithm}
    \end{preview}